\subsection{Finitely generated fields and the subgroup hypothesis}
\paragraph{Statement AE1 (Proved; Part E(a), arithmetic and centralizer assertions).}\label{res:AE1}
The manuscript passages are cote 149, section 3, manuscript pp.~27 and 30, scans 56 and 61--62.
Let $K/\mathbb Q$ be finitely generated, $K'/K$ finite, and $A/K'$ an abelian variety. Then $A(K'(\mu_\infty))_{\rm tors}$ is finite. Every field $L$ finitely generated over a finite extension of $\mathbb Q(\mu_\infty)$, together with every finite extension of $L$, has no nonzero divisible abelian-variety points and has zero invariants in every abelian-variety Tate module. Statements A1--A7, in their indicated centralizer and inertia formulations, consequently extend to finitely generated $K/\mathbb Q$.

\emph{Proof of torsion finiteness.} Replace $K$ by $K'$ and let $k$ be its algebraic constant field over $\mathbb Q$. It is a number field and $K/k$ is regular. Choose a smooth geometrically integral variety $S/k$ with function field $K$, shrink it, and extend $A$ to an abelian scheme $\mathcal A/S$. The extension of the group law, inverse, identity, properness and smoothness uses finite presentation and a common open subset on which the corresponding identities and properties hold. Write $k^{\rm cyc}=k(\mu_\infty)$, so $K(\mu_\infty)=Kk^{\rm cyc}$, and $S_{k^{\rm cyc}}$ is geometrically integral.

For every $n$, $\mathcal A[n]$ is finite etale over $S$, since the characteristic is zero. A $Kk^{\rm cyc}$-rational point of $A[n]$ extends uniquely to a section over the normal scheme $S_{k^{\rm cyc}}$: its closure in the finite etale scheme is a finite birational component, and normality identifies that component with the base. Two such sections agreeing at one geometric point agree everywhere, since their equalizer is both open and closed and the base is connected. Choose a closed point $x$ of $S$ with residue field $k_x$, and one point above it after extending constants. Evaluation gives, compatibly with $n$ and the group laws, an injection
\[
 A(Kk^{\rm cyc})_{\rm tors}\hookrightarrow
 \mathcal A_x(k_x(\mu_\infty))_{\rm tors}.
\]
The target is finite by Ribet's Appendix Theorem 1 with the number field $k_x$. This proves the asserted finiteness without using a specialization theorem that requires finite generation of a Mordell--Weil group.

\emph{Proof about divisible points.} A field $L$ as in the statement is $K_0(\mu_\infty)$ for a finitely generated $K_0/\mathbb Q$: descend its finitely many generators and the finitely many coefficients in their defining relations. The same holds after any finite extension of $L$. Descend an abelian variety and a proposed nonzero divisible point $P$ to $K_0$, then spread them to $\mathcal A/S$ and a section $P$ as above. Shrink $S$ so that $P$ does not meet the identity section. For every $n$ the division scheme $[n]^{-1}(P)$ is finite etale over $S$. A division point over $L$ extends over $S_{k^{\rm cyc}}$ by the same normality argument. Evaluation at $x$ would therefore make the nonzero point $P_x$ divisible by every integer over $k_x(\mu_\infty)$, contradicting A1. Hence there is no such $P$. Invariant Tate vectors vanish either by torsion finiteness or by the compatible nonzero divisible torsion points they would provide.

For completeness, the multiplicative assertion of A1 extends too. The image of the cyclotomic character of $K_0$ is open, since its algebraic constant field is a number field. Its kernel fixes $K_0(\mu_\infty)$, so the same cocycle argument as in A1 shows that an infinitely divisible element has a power with zero Kummer class over $K_0$. The kernel of the latter Kummer map is trivial: on a normal proper model over the algebraic constant field, divisibility by every $n$ forces every divisorial valuation to be zero; a rational function with neither zeros nor poles is algebraic over that constant field. The number-field argument in A1 then applies. Splitting a torus over a finite extension proves the same vanishing of homomorphisms to abelian Tate modules.

Finally every finite geometric covering appearing in A4 descends to a finite extension of this finitely generated $K$, and its compactified Jacobian is an abelian variety there. The torsion finiteness just proved supplies exactly the invariant vanishing used in (A.1). The maps, subgroup containments, compactness and boundary calculations in A2--A7 are unchanged. Thus their proofs transfer, with these specified arithmetic input maps, to this generality.\hfill$\square$

\paragraph{Statement AE2 (Proved; H1 with the necessary zero-centralizer qualification).}\label{res:AE2}
This refines H1 and the subgroup question of cote 149, section 3, manuscript p.~30, scans 61--62.
Let $K$ be a number field, $U/K$ a hyperbolic curve, $s:\Gamma_K\to E_U$ a continuous section, and $\Delta\subset\Gamma_K$ a closed subgroup satisfying
\[
 T(A)^{\Delta\cap\Gamma_{K'}}=0
 \quad\text{for every finite }K'/K\text{ and every abelian variety }A/K'.\tag{E.1}
\]
Write $C_\Delta=C_\pi(s(\Delta))$ and
$P_\Delta=\{\ell:\chi_\ell|_\Delta=1\}$. Every nonzero $C_\Delta$ has a unique associated lifted cusp $c$, fixed by $s(\Delta)$, and
\[
 C_\Delta=I_c^{\chi(\Delta)}
 =\prod_{\ell\in P_\Delta}\mathbb Z_\ell(1).
\]
For $P_\Delta\ne\varnothing$, existence of an $s(\Delta)$-fixed cusp is equivalent to $C_\Delta\ne1$ and that fixed cusp is unique. For $P_\Delta=\varnothing$, $C_\Delta=1$ for every section, including cuspidal sections. When $\Delta$ is normal, existence of an $s(\Delta)$-fixed cusp implies that $s$ itself is cuspidal at that unique cusp, including the case $P_\Delta=\varnothing$.

\emph{Proof.} Use the construction in A4 for $z\in C_\Delta$ and $H=N\langle z\rangle$. The stabilizer of $H$ in $\Gamma_K$ is open; write it as $\Gamma_{K'}$. The image of $z$ in the Tate module of the compactified covering is fixed by $\Delta\cap\Gamma_{K'}$, and (E.1) makes it zero. Thus A3 puts $z$ in inertia, and A5 puts all of $C_\Delta$ in one inertia subgroup whenever $C_\Delta\ne1$. A nonidentity $z\in C_\Delta\cap I_c$ and A2 show that $s(\Delta)$ fixes $c$. The action on $I_c$ is $\chi|_\Delta$, so the fixed subgroup of $I_c$ is exactly $C_\Delta$. Computing each factor in the integral domain $\mathbb Z_\ell$ gives the displayed product. If $P_\Delta$ is nonempty, every fixed cusp contributes a nonzero such product, so A5 and A2 give its uniqueness. If $P_\Delta$ is empty, the same calculation makes a nonzero centralizer impossible.

Suppose now that $\Delta$ is normal and put $L=\overline K^{\Delta}$. We verify the additional assertion without obtaining information from a zero centralizer. Every finite extension $L'/L$, every abelian variety over $L'$ and every one of its rational points descend to some finite $K'/K$ with $L'=LK'$. The group $\Delta\cap\Gamma_{K'}$ is normal in $\Gamma_{K'}$. For a point divisible over $L'$, choose compatible division points there. Its Kummer cocycle is zero on this normal subgroup. Comparing the cocycle on $hg$ and $g(g^{-1}hg)$ puts its value at $g$ in $T(A)^{\Delta\cap\Gamma_{K'}}=0$. The point is therefore divisible over $K'$ and is zero by reduction modulo powers of a good finite prime, as proved in A1. Thus all finite extensions of $L$ have no nonzero divisible abelian-variety points.

The continuous base-change homomorphism $\pi_1(U_L)\to E_U$ identifies its source with the inverse image of $\Delta$: finite etale covers descend to finite intermediate fields, and their subgroup identifications pass to this inverse limit. Thus $s|_\Delta$ is an actual section for $U_L/L$. Stix Theorem 14, with this hypothesis verified as in the second proof of A6, now gives uniqueness of an $s(\Delta)$-fixed lifted cusp. Normality means that $s(\Gamma_K)$ permutes the fixed-cusp set; a singleton is fixed. Hence $s$ is cuspidal at that cusp. Alternatively, for $C_\Delta\ne1$, normality preserves $C_\Delta$ and A2 immediately puts $s(\Gamma_K)$ in $D_c$.\hfill$\square$

The original H1 conclusion that nontrivial centralizers require triviality of the entire cyclotomic character is replaced by the primewise product above. There is no inference from $C_\Delta=1$ to absence of a cusp when $P_\Delta$ is empty. For arbitrary nonnormal $\Delta$, (E.1) alone does not supply the normal-subgroup cocycle argument just given; uniqueness of fixed cusps in this zero-centralizer case is not asserted here. The precise remaining assertion would be that such a subgroup can fix at most one lifted cusp even when all $\chi_\ell|_\Delta$ are nontrivial. Our proof establishes the complete assertion needed for the original, normal-subgroup H1 and for every nonzero centralizer.

\paragraph{Statement AE3 (Proved; arithmetic instances of (E.1), Part E(b)).}\label{res:AE3}
These are the arithmetic instances requested for H1 and cote 149, section 3, manuscript p.~30, scan 62.
For a number field $K$, condition (E.1) holds for the absolute Galois groups of $K(\mu_{\ell^\infty})$, of $K(\mu_\infty)$, and of $K(K^{1/\infty})$. It also holds for the Kummer fields in Murotani--Ozeki Proposition 3.15(2). For $\Delta=\Gamma_{K^{\rm ab}}$, condition (E.1), with all of its quantifiers, holds exactly when $K$ contains no CM subfield.

\emph{Sources and reductions.} The first two cases follow from A1, since $\Gamma_{K(\mu_{\ell^\infty})}\supset\Gamma_{K(\mu_\infty)}$, and intersecting with $\Gamma_{K'}$ preserves that inclusion. For the remaining cases we use the following precise statements from the opened 17 January 2025 manuscript \cite{MurotaniOzeki2025}. Its Proposition 3.15(2) has a characteristic-zero field $k$, a Galois extension $L/k$ containing $\mu_\infty$, open images $\chi_p(\Gamma_k)$ for every prime $p$, and $L$ abelian-variety Kummer-faithful (AVKF); it concludes that every subfield of $L(k^{1/\infty})$ is AVKF. Theorem 3.8, for a number field $K$, equates: absence of a CM subfield of $K$; finiteness of $A(K^{\rm ab})_{\rm tors}$ for every $A/K$; and the AVKF property of $K^{\rm ab}$. AVKF means that every abelian variety over every finite extension has no nonzero divisible rational point.

An AVKF field $M$ has $T_\ell(A)^{\Gamma_{M'}}=0$ for every finite $M'/M$: a nonzero invariant vector would give a nonzero divisible $\ell$-primary torsion subgroup, whose elements are also divisible by integers prime to $\ell$. For an abelian variety defined over $K'$, take $M'=MK'$. This proves (E.1). Apply Proposition 3.15(2) to $k=K$, $L=K(\mu_\infty)$, for which A1 verifies AVKF and the cyclotomic images are open. This gives the asserted full radical field as well as all the stated general instances. No claim of finiteness of the total torsion group over $K(K^{1/\infty})$ is required.

For $K^{\rm ab}$, Theorem 3.8 and the preceding implication prove the forward assertion when $K$ contains no CM subfield. In the opposite case, the proof of that theorem explicitly constructs a finite $K'\subset K^{\rm ab}$ and a nonzero CM abelian variety $A/K'$ all of whose torsion is defined over $K^{\rm ab}$. For a CM subfield $k\subset K$, the inclusion $Kk^{\rm ab}\subset K^{\rm ab}$ follows because base change of an abelian Galois extension remains abelian; base-changing that construction and taking the compositum of its field of definition with $K$ supplies the required $K'$. Thus $T(A)^{\Gamma_{K^{\rm ab}}}=T(A)\ne0$, and (E.1) fails. This also identifies the required finite extension in the failed universal quantifier. The fixed-variety criterion quoted in that source from Zarhin's Theorem 1 concerns $K$-simple factors whose CM endomorphisms are defined over $K$; the quantified assertion (E.1) is the field criterion just proved.\hfill$\square$

\paragraph{Statement AE4 (Refuted original H1; proved counterexample and replacement).}\label{res:AE4}
The original H1 assertion that a nontrivial centralizer requires triviality of the entire cyclotomic character is false. Take $K=\mathbb Q$, $U=\mathbb P^1\setminus\{0,1,\infty\}$, the tangential section at $0$ with tangent vector $1$, and
\[
 \Delta=\operatorname{Gal}(\overline{\mathbb Q}/\mathbb Q(\mu_{3^\infty})).
\]
Then (E.1) holds, $\Delta$ is normal, $\chi|_\Delta$ is nontrivial, and
\[
 C_\pi(s(\Delta))=\mathbb Z_3(1)\subset I_0.
\]
The true statement in place of that H1 assertion is AE2, with its primewise fixed subgroup.

\emph{Proof.} The inclusion $\Gamma_{\mathbb Q(\mu_\infty)}\subset\Delta$ proves (E.1) by A1, after intersection with every $\Gamma_{K'}$. The cyclotomic Galois group over $\mathbb Q$ is the inverse limit of $(\mathbb Z/n\mathbb Z)^*$, hence is $\prod_q\mathbb Z_q^*$ by the Chinese remainder isomorphisms. The image of $\Delta$ consists precisely of tuples whose $3$-component is $1$. Every other prime component remains the full nontrivial $\mathbb Z_q^*$, whereas the $3$-component is trivial. Applying (A.2) therefore gives exactly $\mathbb Z_3(1)$, a nonzero proper subgroup of $I_0$. This verifies all the original hypothesis quantifiers while contradicting its asserted conclusion. The manuscript background is cote 149, section 3, manuscript p.~30, scan 61, together with the supplied H1 text.\hfill$\square$
